\section{Background and Motivation}
\label{sec:intro_background}

\subsection{Cybersecurity and Vulnerability Assessment}

Modern organisations depend on networked software systems for nearly every operational function, and the security of those systems has become a central concern. A security vulnerability is a flaw in software, configuration, or design that an attacker can exploit to gain unauthorised access, execute arbitrary code, disrupt service, or exfiltrate data. Vulnerabilities are catalogued in the Common Vulnerabilities and Exposures (CVE) database after disclosure; the interval between disclosure and remediation deployment frequently extends to months in production environments, leaving systems exposed to any attacker who can reconstruct an exploit from the published description.

\textbf{Vulnerability Assessment and Penetration Testing (VAPT)} is the practice of proactively identifying and demonstrating these weaknesses before an adversary does. A penetration test is an authorised, systematic attempt to discover and exploit weaknesses in a target system, typically following a structured workflow: initial reconnaissance to understand the target's attack surface, identification of candidate weaknesses, exploitation of confirmed vulnerabilities, and reporting of findings with remediation guidance. VAPT is widely recognised as an essential security practice, and demand for it has grown considerably as organisations increase their exposure through cloud services, APIs, and web applications.

\subsection{Why Conventional VAPT Has Limits}

Despite its importance, conventional penetration testing faces practical constraints. Skilled human testers are scarce and expensive: a thorough engagement against a moderately complex target demands days of expert effort and costs can reach thousands of dollars per assessment \citep{kong2025vulnbot}. Scheduled, periodic testing leaves gaps during which newly introduced vulnerabilities remain undetected. Automated scanning tools --- such as network port scanners and web application vulnerability scanners --- can partially fill this gap, but they operate on fixed rule sets and signature databases and lack the contextual reasoning needed to chain together exploits. The practical limits of these tools are not merely theoretical: \citet{fang2024oneday} tested the two most widely deployed open-source scanners (ZAP and Metasploit) against a benchmark of 15 real-world one-day vulnerabilities and found that both achieved 0\% success, whereas a capable LLM agent succeeded on the same targets. Scanners cannot adapt to novel application logic, reason about prerequisite relationships between attack steps, or operate across the breadth of a modern attack surface without human guidance.

These limitations have driven a recent and rapid research interest in applying LLM-based agents to automate parts of the VAPT process --- an interest that has produced a substantive body of work since 2023.

\subsection{Large Language Models and AI Agents}

Large Language Models (LLMs) are deep learning models trained on large corpora of text that have demonstrated a remarkable ability to follow instructions, reason through multi-step problems, write and execute code, and interact with external tools.

An \textbf{LLM agent} extends a base language model with the ability to take actions in an environment: calling tools (such as a web browser, a terminal, or an API), observing the results, and iterating toward a goal. The simplest agents follow a \emph{Reason + Act} (ReAct) loop --- reason about the current state, decide on an action, observe the outcome, and repeat. More sophisticated designs introduce hierarchical planning, role-specialised sub-agents, persistent memory, and structured validation.

A finding that recurs independently across six systems in the VAPT literature --- and that directly shapes the design choices in this thesis --- is that \textbf{architecture, not model scale, is the dominant variable in autonomous exploitation performance} \citep{wang2025surveyagents}. Structured pipelines running comparatively inexpensive models consistently outperform unstructured loops running frontier models: \citet{singer2025incalmo} show that their multi-host framework with Claude Haiku~3.5 surpasses a strong single-agent baseline built on Claude Sonnet~4, and separate work demonstrates that GPT-4o-mini exceeds GPT-4o at web exploitation once a deterministic finite-state machine governs agent behaviour \citep{wang2025surveyagents}. This architectural primacy motivates close attention to structural design rather than model selection in the work described here.

The emergence of capable LLM agents has opened a genuinely new question for cybersecurity: \textit{can an LLM agent carry out penetration testing tasks autonomously, without step-by-step human direction?}

\subsection{The Emergence of Autonomous VAPT Research}

A growing body of research --- developed between 2023 and 2025 --- has answered parts of this question affirmatively, and in doing so has produced the quantitative baseline against which this thesis positions itself.

Early work by \citet{fang2024hackwebsites} showed that a single GPT-4 agent with browser access could autonomously exploit simplified web vulnerabilities, establishing that the basic capability exists. A later study by the same group collected a benchmark of 15 real-world one-day vulnerabilities and found that GPT-4, when given a CVE description, achieved a pass@5 success rate of 87\% --- while every other tested model (GPT-3.5, eight open-source models) and every open-source scanner achieved 0\% \citep{fang2024oneday}. Crucially, without the CVE description, GPT-4's pass@5 rate collapsed to 7\%, establishing an important early distinction: exploiting a vulnerability whose location and nature the agent already knows differs fundamentally from discovering an unknown vulnerability independently. This discovery gap is the central challenge that subsequent systems attempt to close.

\citet{zhu2024hptsa} addressed this gap directly by introducing a hierarchical team of planning and task-specific agents (HPTSA), which was the first system reported to autonomously exploit real zero-day vulnerabilities --- vulnerabilities for which no CVE description is provided. HPTSA achieved a pass@5 rate of 42\% and a pass@1 rate of 18\% on a 14-CVE benchmark, representing a 4.3$\times$ improvement over a no-description single-agent baseline on pass@1. These figures are the primary comparative baseline for the evaluation described in this thesis.

\citet{deng2024pentestgpt} identified two recurring failure modes in single-agent systems: first, context loss over long sessions, which causes the agent to lose awareness of previously attempted attack paths (Finding~3); and second, a depth-first, recency-biased search pattern, in which the agent anchors to the most recently discussed surface and fails to return to earlier discoveries (Finding~4). These two modes together explain why single-agent systems consistently underperform on targets that require coordinated, multi-step attack chains. Later systems address these failure modes through multi-agent architectures \citep{kong2025vulnbot,zhu2024hptsa}, declarative task abstractions \citep{singer2025incalmo}, and planning-augmented designs \citep{wang2025checkmate}.

\subsection{Motivation for This Thesis}

The starting point for this work was a close reading of this body of literature to determine whether it had already addressed the core challenges of autonomous VAPT, and if not, where exactly the remaining gaps lie. The review spans single-agent exploitation studies, hierarchical multi-agent frameworks, planning-augmented approaches, and the principal benchmark suites used to assess autonomous VAPT agents (surveyed in full in Chapter~\ref{chapter_backStudy}).

Two recurring failure modes surfaced across independently developed systems: agents that search broadly over a target's attack surface do so without an explicit model of which vulnerabilities depend on which others; agents that reason about such dependencies do so on pre-enumerated, curated weakness sets that are not assembled dynamically during a live mission. These two failure modes appear to be complementary rather than coincidental, and they are the gaps this thesis has chosen to investigate. The investigation is at an early stage --- the following chapters describe the direction, the proposed design, and the evidence that motivates it, rather than results already obtained.
